% Phase 2 (Guillermo): Part 1 a, b. Every number here comes from src/microstates.py.
\subsection{Enumeration of microstates}
Each microstate of $N$ coins is stored as an integer $s\in\{0,\dots,2^N-1\}$ whose bit $i$ is coin $i$ (1 = heads, 0 = tails). The generator \texttt{microstates(N)} runs over \texttt{range(1 << N)} and yields one state at a time, so memory does not grow with $N$; the macrostate of a state is its number of set bits, $k=$ \texttt{s.bit\_count()}. Counting the states per $k$ gives $\Omega(N,k)$, checked against $\binom{N}{k}$ for $N\le 15$ and against Table~1 of the manual for $N=5$.

\paragraph{(a) $N=7$.} Table~\ref{tab:n7} lists the $2^7=128$ microstates grouped by macrostate. The multiplicities are exact integers (no uncertainty) and follow row 7 of Pascal's triangle: symmetric under heads $\leftrightarrow$ tails and largest at $k=3,4$.

\begin{table}[htbp]\centering\small
\caption{All microstates of $N=7$ coins grouped by macrostate (H = heads, T = tails, coin 0 first).}\label{tab:n7}
\begin{tabular}{@{}p{0.17\textwidth}p{0.68\textwidth}r@{}}
\toprule
Macrostate & Microstates & $\Omega(7,k)$ \\
\midrule
7 heads, 0 tails & \texttt{HHHHHHH} & 1 \\
6 heads, 1 tail & \texttt{HHHHHHT, HHHHHTH, HHHHTHH, HHHTHHH, HHTHHHH, HTHHHHH, THHHHHH} & 7 \\
5 heads, 2 tails & \texttt{HHHHHTT, HHHHTHT, HHHTHHT, HHTHHHT, HTHHHHT, THHHHHT, HHHHTTH, HHHTHTH, HHTHHTH, HTHHHTH, THHHHTH, HHHTTHH, HHTHTHH, HTHHTHH, THHHTHH, HHTTHHH, HTHTHHH, THHTHHH, HTTHHHH, THTHHHH, TTHHHHH} & 21 \\
4 heads, 3 tails & \texttt{HHHHTTT, HHHTHTT, HHTHHTT, HTHHHTT, THHHHTT, HHHTTHT, HHTHTHT, HTHHTHT, THHHTHT, HHTTHHT, HTHTHHT, THHTHHT, HTTHHHT, THTHHHT, TTHHHHT, HHHTTTH, HHTHTTH, HTHHTTH, THHHTTH, HHTTHTH, HTHTHTH, THHTHTH, HTTHHTH, THTHHTH, TTHHHTH, HHTTTHH, HTHTTHH, THHTTHH, HTTHTHH, THTHTHH, TTHHTHH, HTTTHHH, THTTHHH, TTHTHHH, TTTHHHH} & 35 \\
3 heads, 4 tails & \texttt{HHHTTTT, HHTHTTT, HTHHTTT, THHHTTT, HHTTHTT, HTHTHTT, THHTHTT, HTTHHTT, THTHHTT, TTHHHTT, HHTTTHT, HTHTTHT, THHTTHT, HTTHTHT, THTHTHT, TTHHTHT, HTTTHHT, THTTHHT, TTHTHHT, TTTHHHT, HHTTTTH, HTHTTTH, THHTTTH, HTTHTTH, THTHTTH, TTHHTTH, HTTTHTH, THTTHTH, TTHTHTH, TTTHHTH, HTTTTHH, THTTTHH, TTHTTHH, TTTHTHH, TTTTHHH} & 35 \\
2 heads, 5 tails & \texttt{HHTTTTT, HTHTTTT, THHTTTT, HTTHTTT, THTHTTT, TTHHTTT, HTTTHTT, THTTHTT, TTHTHTT, TTTHHTT, HTTTTHT, THTTTHT, TTHTTHT, TTTHTHT, TTTTHHT, HTTTTTH, THTTTTH, TTHTTTH, TTTHTTH, TTTTHTH, TTTTTHH} & 21 \\
1 head, 6 tails & \texttt{HTTTTTT, THTTTTT, TTHTTTT, TTTHTTT, TTTTHTT, TTTTTHT, TTTTTTH} & 7 \\
0 heads, 7 tails & \texttt{TTTTTTT} & 1 \\
\midrule
Total & & 128 \\
\bottomrule
\end{tabular}
\end{table}

\paragraph{(b) Largest $N$ that can be fully enumerated.} We timed the full enumeration and histogram of all $2^N$ states (\texttt{time.perf\_counter}, peak memory with \texttt{tracemalloc}) for growing $N$ on a laptop with an Intel Core i7-1255U (12 threads) and 15.7\,GB of RAM. Two versions were tested, both of which build every microstate and count its heads:
\begin{itemize}
\item the pure-Python generator \texttt{macrostate\_counts}: $750$--$780$\,ns per state, constant memory (no traced allocation);
\item \texttt{macrostate\_counts\_fast}: the same enumeration in blocks of $2^{22}$ consecutive integers (\texttt{numpy.arange}), heads counted with \texttt{numpy.bitwise\_count} and histogrammed with \texttt{numpy.bincount}, the range $0\dots2^N-1$ split among the 12 CPU threads. Memory is fixed by the block size: $71$\,MB per process measured, independent of $N$.
\end{itemize}

\begin{table}[htbp]\centering\small
\caption{Wall-clock time of the full enumeration with \texttt{macrostate\_counts\_fast} (one run per $N$; the time per state varies by about $\pm10\%$ between runs).}\label{tab:benchmark}
\begin{tabular}{@{}rrrr@{}}
\toprule
$N$ & $2^N$ & Time (s) & Time per state (ns) \\
\midrule
30 & $1.07\times10^{9}$ & 3.3 & 3.0 \\
31 & $2.15\times10^{9}$ & 5.5 & 2.6 \\
32 & $4.29\times10^{9}$ & 10.9 & 2.5 \\
33 & $8.59\times10^{9}$ & 23.6 & 2.7 \\
34 & $1.72\times10^{10}$ & 41.4 & 2.4 \\
35 & $3.44\times10^{10}$ & 85.5 & 2.5 \\
36 & $6.87\times10^{10}$ & 156.2 & 2.3 \\
\bottomrule
\end{tabular}
\end{table}

\textbf{Result: $N_{\max}=36$} within 5 minutes ($156$\,s, Table~\ref{tab:benchmark}). Each extra coin doubles the work, so $N=37$ is predicted at $2^{37}\times(2.3\pm0.2)\,\text{ns}\approx310\pm30$\,s, above the limit. With the plain generator the limit is $N=28$ ($2^{28}\times0.76\,\mu\text{s}\approx200$\,s; $N=29$ would take $\approx410$\,s).

\textbf{Limiting factor: time, not memory.} Because the states are generated on the fly and never stored, memory stays constant and the only constraint is the $2^N$ growth of the number of operations: at $2.3$\,ns per state, $N=40$ needs $\approx42$\,min, $N=50$ about a month and $N=100$ about $10^{14}$ years. Memory would become the limit only if the states were kept: a Python list of $2^N$ integers costs $\approx36$\,B per state, which fills the 15.7\,GB of RAM at $N\approx28$, and writing the $2^{36}$ states of $N_{\max}$ as text would take $\approx2.5$\,TB of disk. The fast version gains a factor $\approx330$ over the generator (vectorised popcount in C, $\times100$; 12 threads, $\times3$ since the threads share the memory bandwidth), which is worth only $\log_2 330\approx8$ more coins: exponential growth makes any constant speed-up almost irrelevant.
